\par\noindent\textbf{(1)} 平均场准粒子能量 E(ξ)=√(ξ²+Δ²)。BCS自洽条件为 1=λ∫₀\textasciicircum{}\{ℏω\_D\}[tanh(E/(2k\_BT))/E]dξ，此处假定成对相互作用仅作用于对称Debye窗；E正定，非零Δ是超导破缺相解。
\par\noindent\textbf{(2)} T=0有 tanh(E/2k\_BT)→1，积分为 asinh(ℏω\_D/Δ₀)，故 1/λ=asinh(ℏω\_D/Δ₀)。在Δ₀≪ℏω\_D时 asinh x≈ln(2x)，于是 Δ₀≈2ℏω\_D e\textasciicircum{}\{−1/λ\}；指数非解析于λ，有限阶普通弱耦合微扰无法生成超导能隙。
\par\noindent\textbf{(3)} 在T=T\_c令Δ→0，条件化为 1/λ=∫₀\textasciicircum{}\{ℏω\_D\}tanh[ξ/(2k\_BT\_c)]dξ/ξ。对ℏω\_D≫k\_BT\_c，标准热积分的渐近常数给 k\_BT\_c≈(2e\textasciicircum{}γ/π)ℏω\_D e\textasciicircum{}\{−1/λ\}≈1.134ℏω\_D e\textasciicircum{}\{−1/λ\}，其中γ为Euler常数。
\par\noindent\textbf{(4)} 相除得到 2Δ₀/(k\_BT\_c)=2π/e\textasciicircum{}γ≈3.528，独立于λ和Debye截断，这是弱耦合各向同性BCS的普适比。强耦合、各向异性节点、频率依赖相互作用或明显变化的态密度都可能改变比值，不能把3.528当成所有超导材料的严格常数。
